\begin{ExerciseBlock}
\ExerciseHeading{10}{Python: Bewegungsdiagramme}{linear_kinematics_python_motion_diagrams}
Mit dem Code-Beispiel lassen sich aus einer vorgegebenen Ortsfunktion $s(t)$ die Geschwindigkeitsfunktion $v(t)$ und die Beschleunigung $a(t)$ berechnen.

\begin{enumerate}
	\item Versuchen Sie den Ablauf des Programms zu verstehen, sodass Sie diesen erklären und anpassen können.

	\item Wählen Sie 3 verschiedene Ortsfunktionen $s(t)$ und überlegen Sie sich den groben Verlauf der $v(t)$ und $a(t)$ Funktionen. Überprüfen Sie mit dem Code anschließend ihre Vermutung.

\end{enumerate}

Verwenden Sie das folgende Codebeispiel als Ausgangspunkt.
\PythonCode{content/resources/chapters/linear_kinematics/python/motion_diagrams.py}
\end{ExerciseBlock}
\UseExerciseSolution{linear_kinematics_python_motion_diagrams}
